% ---------------------------------------------------------------------------
% fragment.tex — Gem geometry of the 880 n=4 squares, exact swap-gap
% verification, Siamese gems.  Written by the gem_geometry subagent for v5.
% Assumes amsthm environments proposition/lemma/corollary/remark and \Cov.
% Figures: Analysis/revision/gem_geometry/fig_gem_geometry_880.pdf,
%          Analysis/revision/gem_geometry/fig_gems_cross_order.pdf
% Numbers: Analysis/revision/gem_geometry/results.json
% ---------------------------------------------------------------------------

% ===========================================================================
\subsection{Swap gaps: exact verification}
\label{sec:swap-gaps}

Throughout, $N=n^2$ and $E_{\mathrm{full}}$ is the symmetric energy
$E_{\mathrm{full}}(S)=\sum_{L}\Cov(\mathbf 1_L,Z)^2$, the sum running over all
$2n+2$ lines $L$ ($n$ rows, $n$ columns, two diagonals), with
$\Cov(\varphi,Z)=\frac1N\sum_{ij}(\varphi_{ij}-\bar\varphi)z_{ij}$.
For two distinct cells $a,b$ let $m(a,b)$ be the number of lines containing
exactly one of $a,b$; since two distinct cells differ in their row or in their
column, $m(a,b)\in\{2,\dots,6\}$.

\begin{proposition}[single-swap gap; verification]
\label{prop:swap-verified}
Let $S$ be a magic square of order $n$ and let $S'$ be obtained by exchanging
the entries of two cells $a,b$ with $d=|s_a-s_b|$.  Then
\[
E_{\mathrm{full}}(S')-E_{\mathrm{full}}(S)=\frac{d^{2}\,m(a,b)}{N^{2}} .
\]
The formula was verified in exact rational arithmetic on every one of the
$880\cdot\binom{16}{2}=105{,}600$ single swaps of the $880$ essentially
different order-$4$ squares and the $8\cdot\binom 92=288$ swaps of the eight
order-$3$ squares, and on all swaps of the Siamese squares of orders $5,7,9$:
agreement in $100\%$ of cases, no exceptions.
\end{proposition}

(The proof is in Section~\ref{sec:swap-formula}; here we only record the
computation.)  Two consequences are immediate.  First, every magic square is a
strict local minimum of $E_{\mathrm{full}}$ under transpositions, for every
$n$, since $d\ge 1$ and $m\ge 2$.  Second, the gap spectrum is
\emph{arithmetic}, not ``banded'': over the $105{,}600$ order-$4$ swaps the
quantity $N^2\Delta E=d^2m$ takes exactly $70$ distinct integer values from $2$
to $980$, and the three smallest, $2,3,4$, are precisely the three values that
can occur as the \emph{minimum} gap of a square
(Table~\ref{tab:min-gaps}).  The ``trimodal band structure'' reported in v4
was the histogram of $d^2m/N^2$ computed with a non-symmetric energy; it
carries no information beyond Proposition~\ref{prop:swap-verified}.

\begin{table}[htbp]
\centering
\caption{Minimum single-swap gap $\Delta_{\min}(S)=\min_{a\neq b}\Delta E$
(exact) over the magic squares of orders $3$ and $4$ and the Siamese squares
of orders $5$--$9$.  All values are $k/N^2$ with $k\in\{2,3,4\}$, as forced by
$\Delta E=d^2m/N^2$ with $d=1$.}
\label{tab:min-gaps}
\begin{tabular}{lccl}
\toprule
squares & $N^2$ & $N^2\Delta_{\min}$ & multiplicity \\
\midrule
all $8$ of order $3$ & $81$ & $3$ & $8$ \\
all $880$ of order $4$ & $256$ & $2$ / $3$ / $4$ & $432$ / $420$ / $28$ \\
Siamese, $n=5$ & $625$ & $3$ & -- \\
Siamese, $n=7$ & $2401$ & $2$ & -- \\
Siamese, $n=9$ & $6561$ & $2$ & -- \\
\bottomrule
\end{tabular}
\end{table}

% ===========================================================================
\subsection{Geometry of the 880 order-4 gems}
\label{sec:gem-geometry-880}

For each of the $880$ essentially different order-$4$ magic squares we
computed the Magic Gem $G(S)=\operatorname{conv}\{p_{ij}\}$ exactly: doubling
all coordinates makes every point an integer vector, facet planes are
recovered as primitive integer equations (so coplanar triangles returned by
\textsc{qhull} are merged into true faces), vertices are the points lying on
at least three facet planes, edges are pairs of vertices sharing two facets,
and volumes are computed as exact rationals.  Euler's relation $V-E+F=2$ holds
for all $880$ gems.  Figure~\ref{fig:gem-geometry-880} and
Table~\ref{tab:gem-fvectors} summarise the results.

\begin{table}[htbp]
\centering
\caption{Face structure of the $880$ order-$4$ gems by number of hull
vertices $V$.  ``types'' is the number of distinct combinatorial types (face
lattices up to isomorphism) and ``vols'' the number of distinct exact
volumes among the squares with that $V$.  The $f$-vector column lists
$(V,E,F)$ with multiplicities.}
\label{tab:gem-fvectors}
\footnotesize
\begin{tabular}{rrrrp{7.6cm}}
\toprule
$V$ & squares & types & vols & $f$-vectors $(V,E,F)$: count \\
\midrule
8 & 7 & 4 & 4 & (8,14,8):1, (8,16,10):1, (8,17,11):1, (8,18,12):4 \\
9 & 10 & 1 & 5 & (9,20,13):10 \\
10 & 103 & 27 & 51 & (10,19,11):1, (10,20,12):3, (10,21,13):11, (10,22,14):75, (10,23,15):11, (10,24,16):2 \\
11 & 80 & 32 & 33 & (11,22,13):2, (11,23,14):6, (11,24,15):28, (11,25,16):38, (11,26,17):6 \\
12 & 380 & 129 & 111 & (12,22,12):1, (12,24,14):37, (12,25,15):25, (12,26,16):164, (12,27,17):75, (12,28,18):77, (12,29,19):1 \\
13 & 106 & 52 & 40 & (13,27,16):4, (13,28,17):12, (13,29,18):46, (13,30,19):42, (13,31,20):2 \\
14 & 162 & 85 & 65 & (14,28,16):1, (14,29,17):2, (14,30,18):36, (14,31,19):17, (14,32,20):106 \\
15 & 14 & 7 & 6 & (15,32,19):4, (15,33,20):6, (15,34,21):4 \\
16 & 18 & 15 & 14 & (16,32,18):1, (16,34,20):8, (16,36,22):9 \\
\midrule
all & 880 & 352 & 155 & 39 distinct $f$-vectors \\
\bottomrule
\end{tabular}
\end{table}

\paragraph{Vertex counts, $f$-vectors, volumes.}
Between $8$ and all $16$ of the points are hull vertices; $V=12$ is modal
($380/880$).  There are $39$ distinct $f$-vectors, and the multiset of facet
sizes (number of triangles, quadrilaterals, \dots) is determined by the
$f$-vector in every case.  The exact volume takes $155$ distinct values, from
$58$ (square \#427) to $292/3$ (square \#379), mean $84.05$; denominators are
$1,2,3,6$ as they must be for lattice polytopes of the doubled lattice.
Surface areas range from $116.5$ to $148.6$.

\paragraph{Combinatorial types.}
Two gems have the same combinatorial type if their vertex--facet incidence
structures (equivalently, their face lattices) are isomorphic.  Grouping by a
Weisfeiler--Lehman hash of the bipartite vertex--facet graph and confirming
every group by an explicit graph isomorphism test gives
\[
\textbf{352 combinatorial types}
\]
among the $880$ gems (the hash was already exact: no hash class had to be
split).  Their sizes are very uneven: $93$ types occur once, $201$ twice,
$8$ three times, $20$ four times, and the remaining $30$ types have between
$5$ and $32$ members; the largest type, with $f=(10,22,14)$ (12 triangles,
2 quadrilaterals), is shared by $32$ squares, and the second largest,
$f=(12,24,14)$, by $26$.  The $7$ gems with $8$ vertices fall into $4$ types,
including a single $(8,14,8)$ gem with four triangular and four quadrilateral
faces (square \#54, Fig.~\ref{fig:gem-geometry-880}d) and four
$(8,18,12)$ gems sharing one simplicial type with $12$ triangular faces.  The $18$ gems with all $16$ points
extreme belong to $15$ types.

\paragraph{Congruence classes and injectivity.}
Two gems are \emph{congruent} if an isometry of $\mathbb R^3$ (rotations,
reflections, translations) maps one onto the other.  Because a
distance-preserving bijection between finite point sets extends to an
isometry, congruence of gems is decided exactly by testing for a bijection of
hull-vertex sets preserving all pairwise (integer) squared distances; we
grouped candidates by the sorted multiset of squared distances and confirmed
each group by an edge-labelled graph isomorphism, with a floating-point
Kabsch superposition as a cross-check (maximal RMSD $2\cdot10^{-14}$).

\begin{proposition}[class $\to$ gem]
\label{prop:gem-injective}
\begin{enumerate}
\item The $880$ hull-vertex point sets are pairwise distinct: the map from
  $D_4$-classes of order-$4$ magic squares to gems (as subsets of
  $\mathbb R^3$) is injective.
\item Up to congruence there are exactly $616$ gems.  The congruence classes
  are precisely the orbits of the complement map $S\mapsto 17-S$, which acts
  on gems as the reflection $z\mapsto -z$: $352$ squares are
  self-complementary (their gems are mirror-symmetric up to a lattice
  symmetry) and the other $528$ form $264$ mirror pairs.  No two gems of
  non-complementary squares are congruent, and the same holds for the full
  $16$-point configurations.
\end{enumerate}
\end{proposition}

Part (2) is the only failure of injectivity, and it is the expected one:
$17-S$ is magic whenever $S$ is, and the gem of $17-S$ is the mirror image of
$G(S)$ in the plane $z=0$.  The $352$ self-complementary classes are exactly
Dudeney's group III (the $48$ associated squares, whose complement is the
$180^\circ$ rotation) together with all $304$ squares of Dudeney's
group VI; the complement map preserves every Dudeney group.  The coincidence
$352 = \#\{\text{combinatorial types}\} = \#\{\text{self-complementary
squares}\}$ is numerical only: the type partition ($93$ singletons) and the
complement partition ($352$ singletons) are different partitions of the $880$.

\paragraph{Dudeney groups.}
Table~\ref{tab:dudeney-gems} crosses the hull-vertex count with Dudeney's
twelve groups.  Geometry is only loosely tied to Dudeney type: every group
with more than eight members spans at least four vertex counts and at least
$20$ combinatorial types, while the two $8$-member groups XI and XII are
confined to $V\in\{12,13\}$ and $V\in\{10,12,13\}$ with $4$ types each.  The
$8$-vertex gems occur only in groups VI (three) and VIII (four), and
$V=9$ occurs only in groups VIII and IX.

\begin{table}[htbp]
\centering
\caption{Hull-vertex count $V$ versus Dudeney group for the $880$ order-$4$
squares; last column: number of combinatorial types in the group.}
\label{tab:dudeney-gems}
\small
\begin{tabular}{rr|rrrrrrrrr|r}
\toprule
group & size & $V{=}8$ & 9 & 10 & 11 & 12 & 13 & 14 & 15 & 16 & types \\
\midrule
I & 48 &  &  &  & 10 & 20 & 12 & 4 & 2 &  & 20 \\
II & 48 &  &  & 4 & 14 & 14 & 6 & 6 & 2 & 2 & 22 \\
III & 48 &  &  & 13 &  & 28 &  & 4 &  & 3 & 30 \\
IV & 96 &  &  & 12 & 20 & 40 & 14 & 10 &  &  & 43 \\
V & 96 &  &  & 12 & 10 & 34 & 20 & 16 & 4 &  & 48 \\
VI & 304 & 3 &  & 40 &  & 150 &  & 98 &  & 13 & 126 \\
VII & 56 &  &  &  & 4 & 20 & 22 & 4 & 6 &  & 28 \\
VIII & 56 & 4 & 8 & 10 & 4 & 24 & 6 &  &  &  & 20 \\
IX & 56 &  & 2 & 10 & 18 & 12 & 10 & 4 &  &  & 25 \\
X & 56 &  &  &  &  & 32 & 8 & 16 &  &  & 26 \\
XI & 8 &  &  &  &  & 4 & 4 &  &  &  & 4 \\
XII & 8 &  &  & 2 &  & 2 & 4 &  &  &  & 4 \\
\bottomrule
\end{tabular}
\end{table}

\begin{remark}
Groups III and VI, the self-complementary ones, have only even vertex counts;
this is forced: for a self-complementary square the reflection $z\mapsto-z$
composed with the element of $D_4$ that realises $17-S$ is an isometry of
the gem of order $2$ or $4$ without fixed hull vertices (no point has $z=0$
when $N$ is even), so hull vertices fall into orbits of size $2$ or $4$.
\end{remark}

\begin{figure}[htbp]
\centering
\includegraphics[width=\textwidth]{figures/fig_gem_geometry_880.pdf}
\caption{Geometry of the $880$ order-$4$ Magic Gems.
(a)~Number of hull vertices ($8$ to $16$; modal value $12$).
(b)~Exact hull volume ($155$ distinct values between $58$ and $292/3$).
(c)~Sizes of the $352$ combinatorial types: $93$ types occur once, $201$
twice, and one type is shared by $32$ squares.
(d)--(f)~Three extreme gems, drawn with translucent faces after merging
coplanar triangles: the unique $(8,14,8)$ gem (square \#54; the eight
interior points are shown in red), a gem with all $16$ points extreme
(square \#802, a row-permuted D\"urer square), and the gem of largest
volume (square \#379).}
\label{fig:gem-geometry-880}
\end{figure}

% ===========================================================================
\subsection{Siamese gems have twelve vertices}
\label{sec:siamese-gems}

For the Siamese (de la Loub\`ere) squares of odd order the hull-vertex count
does not grow with $n$.  We computed the gems for $n=3,5,\dots,21$: the Lo Shu
gem has $f=(8,16,10)$, and every Siamese gem with $n\ge5$ has
$f=(12,26,16)$ with $12$ triangular and $4$ quadrilateral facets, so the hull
fraction $12/n^2$ tends to zero (Fig.~\ref{fig:gems-cross-order}).  This is a
theorem, with a short computer-assisted proof.

\begin{proposition}[Siamese gems]
\label{prop:siamese-12}
Let $n=2m+1\ge5$ and let $S$ be the Siamese square of order $n$ ($1$ in the
middle cell of the top row, successive entries one step up and to the right
modulo $n$, one step down when blocked).  Then $G(S)$ has exactly $12$
vertices, $26$ edges and $16$ facets ($12$ triangles, $4$ quadrilaterals).
With $0$-indexed cells $(i,j)$ and $c=(n^2+1)/2$, the vertices are the
points over the cells
\[
\begin{array}{lll}
(0,0)\colon c+m+2, & (0,m-1)\colon n^2-1, & (0,m)\colon 1, \\
(0,2m)\colon c+m, & (m-1,0)\colon n^2-m, & (m,0)\colon m+2, \\
(m,2m)\colon n^2-m-1, & (m+1,2m)\colon m+1, & (2m,0)\colon c-m, \\
(2m,m)\colon n^2, & (2m,m+1)\colon 2, & (2m,2m)\colon c-m-2,
\end{array}
\]
i.e.\ the four corners, and on each side of the square the two cells holding
the side's extreme values.  For $n=3$ the same construction gives the Lo Shu
gem, with $8$ vertices (all cells except the centre).
\end{proposition}

\begin{proof}
\emph{Step 1: closed form.}  Write $\langle t\rangle\in\{0,\dots,n-1\}$ for
$t\bmod n$.  We claim $s(i,j)=n\langle i+j-m\rangle+\langle i+2j-2m\rangle+1$.
Indeed, along a run of $n$ consecutive entries the move $(i,j)\to(i-1,j+1)$
fixes $i+j$ and increases $i+2j$ by $1$ modulo $n$, while the ``blocked''
move $(i,j)\to(i+1,j)$ increases both $i+j$ and $i+2j$ by $1$; the entry $1$
sits at $(0,m)$ where both expressions vanish.  Hence by induction the entry
$k$ lies in a cell with $\langle i+j-m\rangle=\lfloor (k-1)/n\rfloor$ and
$\langle i+2j-2m\rangle=\langle k-1\rangle$; since
$(i,j)\mapsto(i+j,\,i+2j)$ is a bijection of $\mathbb Z_n^2$ (determinant
$1$), this determines the cell, which proves the claim.

\emph{Step 2: affine reduction to five layers.}  In gem coordinates
$x=j-m$, $y=m-i$ put $u=i+j-m=x-y+m$ and $w=i+2j-2m=2x-y+m$.  Since
$\langle t\rangle=t-n\lfloor t/n\rfloor$,
\[
z=s-\tfrac{n^2+1}{2}=\underbrace{nu+w+1-\tfrac{n^2+1}{2}}_{L(x,y)\ \text{affine}}
+\;n\,\zeta(x,y),\qquad
\zeta:=-n\Big\lfloor\frac{x-y+m}{n}\Big\rfloor-\Big\lfloor\frac{2x-y+m}{n}\Big\rfloor .
\]
The map $\Phi(x,y,z)=(x,y,(z-L(x,y))/n)$ is an affine bijection of
$\mathbb R^3$, so $\Phi(G(S))=\operatorname{conv}\{(x,y,\zeta(x,y))\}$ has the
same face lattice as $G(S)$ and the same vertex cells.  Because
$-m\le x-y+m\le 3m$ and $-2m\le 2x-y+m\le 4m$, both floors lie in
$\{-1,0,1\}$; the floor of $(x-y+m)/n$ equals $-1$ iff $y\ge x+m+1$ and
$1$ iff $y\le x-m-1$, that of $(2x-y+m)/n$ equals $-1$ iff $y\ge 2x+m+1$ and
$1$ iff $y\le 2x-m-1$.  The combinations $(-1,1)$ and $(-1,0)$ are impossible
in the box $|x|,|y|\le m$ (they force $x\ge 2m+2$, resp.\ $x\ge1$ together
with $y>m$), and by the central symmetry $\zeta(-x,-y)=-\zeta(x,y)$ so are
$(1,-1),(1,0)$.  The configuration therefore consists of the lattice points of
five convex lattice polygons at five heights:
\[
\begin{array}{ll}
T=\operatorname{conv}\{(-m,m),(-1,m),(-m,1)\}, & \zeta=n+1;\\
U=\operatorname{conv}\{(-m,1-m),(-m,0),(-1,m-1)\}, & \zeta=1;\\
M=\operatorname{conv}\{(-m,-m),(0,m),(m,m),(0,-m)\}, & \zeta=0;\\
D=-U,\quad B=-T, & \zeta=-1,\ \zeta=-(n+1).
\end{array}
\]
(For $T$: $y\ge x+m+1$ with $y\le m$ forces $x\le-1$ and then implies
$y\ge 2x+m+1$; for $U$: $2x+m+1\le y\le x+m$ forces $x\le -1$ and implies the
other constraints; $M$ is cut out by $y\le 2x+m$, $y\ge 2x-m$, $|y|\le m$.)

\emph{Step 3: the twelve points and sixteen planes.}  Let
$v_1=(-m,m,n{+}1)$, $v_2=(-1,m,n{+}1)$, $v_3=(-m,1,n{+}1)$, $v_4=(-m,0,1)$,
$v_5=(-m,-m,0)$, $v_6=(0,-m,0)$, $v_7=(0,m,0)$, $v_8=(m,m,0)$, $v_9=(m,0,-1)$,
$v_{10}=-v_1$, $v_{11}=-v_3$, $v_{12}=-v_2$ (so $v_{10}=(m,-m,-n-1)$,
$v_{11}=(m,-1,-n-1)$, $v_{12}=(1,-m,-n-1)$; these are the cells listed in
the statement).  Consider the eight inequalities
\begin{align*}
&(1)\ \zeta\le 2m+2, \qquad (2)\ y\le m, \qquad (3)\ x\ge -m, \qquad (8)\ -2y+\zeta\le 2m,\\
&(4)\ -2(m{+}1)x+(2m{+}1)y-m\zeta\le m(2m{+}1),\\
&(5)\ 2(m{+}1)x-2(m{+}1)y+(2m{+}1)\zeta\le 2m(m{+}1),\\
&(6)\ 2m(m{+}2)x-(2m^2{+}2m{-}1)y+m(2m{+}1)\zeta\le m(2m^2{+}2m{-}1),\\
&(7)\ 2mx-y+m\zeta\le m(2m{-}1),
\end{align*}
together with their images $(1')$--$(8')$ under $(x,y,\zeta)\mapsto
(-x,-y,-\zeta)$, and let $Q$ be the polyhedron they define.  The following
facts are finitely many polynomial identities and inequalities in $m$; they
were verified exactly (integer polynomial arithmetic; an inequality
$p(m)>0$ for all integers $m\ge2$ is certified by writing $p(2+t)$ with
non-negative coefficients, or, failing that, by checking $p$ exactly on
$[2,R]$ for a Cauchy root bound $R$ and the sign of the leading coefficient
beyond) by the script \texttt{siamese\_proof.py}:
\begin{enumerate}
\item[(a)] every vertex of the five polygons $T,U,M,D,B$, placed at its
  height, satisfies all $16$ inequalities; by linearity so does every
  configuration point, i.e.\ $\Phi(G(S))\subseteq Q$;
\item[(b)] each plane contains exactly the following $v$'s, with every other
  $v_k$ strictly inside (and the centrally symmetric lists for
  $(1')$--$(8')$):
  \[
  \begin{array}{llll}
  (1)\colon v_1v_2v_3, & (2)\colon v_1v_2v_8v_7, & (3)\colon v_1v_3v_5v_4, & (4)\colon v_1v_4v_7,\\
  (5)\colon v_2v_3v_6, & (6)\colon v_2v_6v_9, & (7)\colon v_2v_8v_9, & (8)\colon v_3v_5v_6;
  \end{array}
  \]
\item[(c)] the two quadrilaterals (and their mirrors) are in convex position
  in the cyclic order given, so the $16$ face polygons have $26$ distinct
  edges, each shared by exactly two polygons, and $12-26+16=2$;
\item[(d)] each $v_k$ lies on three of the planes with linearly independent
  normals.
\end{enumerate}
By (b), each of the $16$ planes supports $P:=\operatorname{conv}\{v_1,\dots,
v_{12}\}$ in a $2$-dimensional face, hence a facet, of $P$ with the stated
vertices.  By (c) this family of facets is closed under edge-adjacency; the
facet-adjacency graph of a $3$-polytope is connected, so these are
\emph{all} facets of $P$ and $P=Q$.  With (a),
$P\subseteq\Phi(G(S))\subseteq Q=P$, so $\Phi(G(S))=P$, whose vertex set is
$\{v_1,\dots,v_{12}\}$ by (d).  Mapping back by $\Phi^{-1}$ gives the
statement; the facet counts are read off from (b)--(c).
\end{proof}

\begin{remark}
For $m=1$ ($n=3$) the triangles $T$, $U$, $D$, $B$ degenerate to single
points, $v_1=v_2=v_3$ and $v_{10}=v_{11}=v_{12}$, and the same argument gives
the $8$-vertex Lo Shu gem.  The proof shows more than the count: up to an
affine map, every Siamese gem is the convex hull of five parallel lattice
polygons whose shapes do not depend on $n$ except through scale, which is why
the combinatorial type is constant for all $n\ge5$.  Whether other classical
constructions (e.g.\ the LUX or Strachey methods for $n\equiv 0,2\pmod 4$)
have bounded hull-vertex counts is open.
\end{remark}

% ===========================================================================
\subsection{Cross-order comparison}
\label{sec:cross-order}

\begin{figure}[htbp]
\centering
\includegraphics[width=\textwidth]{figures/fig_gems_cross_order.pdf}
\caption{Magic Gems across orders, drawn with coplanar triangles merged into
true faces (translucent) and interior points in red: (a)~the Lo Shu
($n=3$, $8$ of $9$ points extreme), (b)~D\"urer's square ($n=4$, $12$ of
$16$), (c)~the Siamese square of order $5$ ($12$ of $25$;
Proposition~\ref{prop:siamese-12}).  Exact volumes $18$, $92$, $830/3$.}
\label{fig:gems-cross-order}
\end{figure}
